\documentclass{article}
\usepackage[T1]{fontenc}
\usepackage{lmodern}
\usepackage{graphicx}
\usepackage{amsmath,amssymb}
\usepackage[a4paper,margin=2cm]{geometry}
\usepackage[absolute,overlay]{textpos}
\setlength{\TPHorizModule}{1mm}
\setlength{\TPVertModule}{1mm}
\usepackage[indonesian]{babel}
\usepackage{hyperref}
\setlength{\emergencystretch}{2em}
% unit: Halaman 28

\begin{document}

% === Halaman 28 (llm) ===

\textcolor{red}{(02 \bullet 26) \oplus (03 \bullet 7B) \oplus (01 \bullet BD) \oplus (01 \bullet 43) = 3F}

Selanjutnya, XOR-kan semua hasil antara tersebut:

\begin{align*}
(02 \bullet 26) &= 4C = 0100 \ 1100 \\
(03 \bullet 7B) &= 8D = 1000 \ 1101 \\
(01 \bullet BD) &= BD = 1011 \ 1101 \\
(01 \bullet 43) &= 43 = \underline{0100 \ 0011} \oplus \\
& \quad 0011 \ 1111 = 3F
\end{align*}

Jadi, $(02 \bullet 26) \oplus (03 \bullet 7B) \oplus (01 \bullet BD) \oplus (01 \bullet 43) = 3F$
Persamaan lainnya diselesaikan dengan cara yang sama.

\end{document}
